% Generated draft: requires mathematical and source review.
\documentclass[11pt]{article}
\usepackage[margin=25mm]{geometry}
\usepackage{fontspec}
\setmainfont{Latin Modern Roman}
\usepackage{amsmath,amssymb,mathtools,unicode-math}
\setmathfont{Latin Modern Math}
\usepackage{enumitem,needspace,xurl,hyperref}
\setlength{\parindent}{0pt}
\setlength{\parskip}{5pt}
\setlength{\emergencystretch}{3em}
\newcommand{\N}{\mathbb{N}}
\newcommand{\Z}{\mathbb{Z}}
\newcommand{\Q}{\mathbb{Q}}
\newcommand{\R}{\mathbb{R}}
\newcommand{\C}{\mathbb{C}}
\newcommand{\F}{\mathbb{F}}
\newcommand{\A}{\mathbb{A}}
\newcommand{\PP}{\mathbb{P}}
\newcommand{\E}{\mathbb{E}}
\newcommand{\eps}{\varepsilon}
\newcommand{\dd}{\,\mathrm{d}}
\begin{document}
\section*{UnsolvedMath ID: 20000874}
\textbf{Source problem code:} catalogue entry
\subsection{Definitions and mathematical statement}
This section is a moderated discussion of metric choice on tree space, statistical continuity, intrinsic versus extrinsic geometry, likelihood and residuals, and communication with biologists. It contains several questions rather than one conjecture. The supplied theorem addresses its Robinson--Foulds/BHV comparison and one family of nonpositively curved metrics; it does not choose a universally optimal biological geometry.

Fix a labeled leaf set with at least four leaves and the usual unrooted BHV tree space, formed by gluing Euclidean nonnegative orthants of compatible internal-edge splits along zero-length faces. Pendant-edge lengths are omitted or held fixed. For trees \(T,T'\), let \(S,S'\) be their positive internal split sets, with lengths \(\ell_s,\ell'_s\). Define the unnormalized topology-only Robinson--Foulds distance and common-split discrepancy by
\[
 r=d_{\mathrm{RF}}(T,T')=|S\mathbin{\triangle}S'|,\qquad
 D_C^2=\sum_{s\in S\cap S'}(\ell_s-\ell'_s)^2.
\]
On metric trees the topology-only quantity is a pseudometric, since changing shared lengths can leave \(r=0\). If variable pendant lengths are included, their common-coordinate discrepancy must also be included in \(D_C\).

Assume every unmatched positive length lies in \([\delta,M]\), \(0<\delta\leq M\). The aid's sharp certificate is
\[
 \sqrt{D_C^2+\delta^2r}\ \leq\ d_{\mathrm{BHV}}(T,T')
 \ \leq\ \sqrt{D_C^2+2M^2r}.
\]
The lower unmatched-length coefficient is sharp for a star/refinement comparison; the common-length term is sharp in one orthant; and the upper factor two is sharp for two incompatible quartet splits of equal length \(M\), whose distance is \(2M\) and whose RF count is two. Without a positive lower unmatched-length bound, fixed topological differences can occur at arbitrarily small BHV distance. Without tracking common lengths, RF zero can coexist with positive BHV distance. Even with \(D_C=0\), no upper control by RF is possible without an upper unmatched-length bound, since lengths can be scaled arbitrarily. Thus a purported unconditional Lipschitz comparison must specify its direction, normalization, scale restrictions, and treatment of shared lengths.

For the curvature question, assign each possible split a positive weight \(w_s\), using the same weight whenever that split appears in an orthant. Give each orthant the diagonal length metric \(\sum_s w_s\,\dd\ell_s^2\) and glue consistently. The map \(\ell_s\mapsto\sqrt{w_s}\ell_s\) is a global isometry to ordinary BHV space, so the weighted metric is CAT(0). Nonconstant weights are not scalar multiples of the ordinary metric on the fixed physical-length coordinates. This gives coordinate-dependent nonscalar examples, not a classification of all CAT(0) geometries up to arbitrary isometry or under additional biological symmetry restrictions.

The discussion's other qualifications remain relevant. St. John and Penny discussed \(L^1\) distances and time scaling, while Felsenstein was associated with \(L^2\); Flath and Penny raised differing importance of deep internal versus near-leaf edges, and topology-only versus length-sensitive tasks. Forman wanted continuity of statistical methods and attention to residuals; Charney, Vogtmann and others questioned which ambient space or intrinsic metric was appropriate. The sharp comparison does not resolve those choices.

Vert's likelihood discussion fixes an evolutionary probability model taking a tree to probabilities on leaf-letter assignments in \(\{A,G,C,T\}^{L}\). With empirical distribution \(\widehat p\) and model distribution \(p_T\), maximizing log likelihood corresponds to minimizing \(D_{\mathrm{KL}}(\widehat p\Vert p_T)\), with the usual support conventions. KL is generally asymmetric and may be infinite, and a model map is not automatically an injective embedding. The report's warning that an intrinsic average of high-likelihood trees need not retain high likelihood is not resolved by an RF/BHV bound. No substitution model, ambient divergence, averaging rule, or statistical optimality theorem is supplied.

The remaining discussion proposed clearer review articles, possible Science outreach, a venue called ``Systematics Biology'' in the notes, cross-disciplinary courses, and explaining useful notions of average. These are historical communication suggestions, not mathematical assertions or instructions to contact or publish anything.
\subsection{Short English statement}
The supplied theorem makes RF-versus-BHV comparison precise by retaining unmatched-edge scale bounds and shared-length differences, with sharp constants. Split-consistent positive weights give further CAT(0) metrics by rescaling coordinates, though they remain isometric to BHV. The broader choices of biological metric, likelihood geometry, averaging, and communication remain discussion questions.
\subsection{Sources}
\begin{itemize}
\item Local registry: permanent UnsolvedMath ID \(20000874\), exact public code \texttt{catalogue entry}; \texttt{aim-biology-notes.json}, \texttt{notes[38]}. AIM domain Biology; registry category Miscellaneous; source tag section.
\item AIM notes \emph{Geometric Models of Biological Phenomena}, version dated 18 June 2003, Section A.5, ``Wednesday Discussion on Geometric Issues,'' moderated by John Shareshian: \url{https://aimath.org/WWN/geombio/geombio.pdf}. This URL is supplied provenance; no PDF access is claimed.
\item Recorded participants include Vogtmann, Forman, Bridson, Charney, Billera, St. John, Penny, Flath, Diaconis, Vert, Snel, and Huelsenbeck; Felsenstein's approaches are discussed. The source attributes the RF/BHV comparison question to Diaconis and the Lipschitz-constant question to Bridson.
\item Supplied research summary, ``Sharp scale-bounded comparison of BHV and Robinson-Foulds distance,'' attempt 1. The discussion names BHV, Robinson--Foulds and relative entropy but supplies no complete mathematical references or empirical model citations; original remarks and literature fields are empty.
\end{itemize}
\end{document}
